\section{The full iPRA interpretation, explicitly}
\label{sec:interpretations}

The exact consequence theorem in Section~\ref{sec:fullnf-metatheory}
reduces the arithmetic-strength question to the raw native calculus
$\mathcal N$. This section states the connecting interpretations and
the inherited arithmetic dependency order explicitly. The appendix gives
the underlying proofs, rather than replacing them with standard-model
calculations.

\subsection{Recovering ordinary intuitionistic logic}

In raw scalar syntax define the numerical selector
\[
 C(v,a,b)=I(v,a;u.b),
 \qquad C(0,a,b)=a,\quad C(Sv,a,b)=b.
\]
For the normalized presentation take the normal form of each displayed
term. A formula gate $G_v A$ is obtained by replacing every equation
$s=t$ with $C(v,0,s)=C(v,0,t)$ and passing the gate through $Q$ and $D$,
renaming bound variables to avoid capture.
It satisfies
\[
 \vdash G_0A,\qquad G_{Sv}A\dashv\vdash A.
\]
The case rules and binder congruence extend a numerical selector to a
formula selector $M_v(A,B)$ with
\[
 M_0(A,B)\dashv\vdash A,\qquad M_{Sv}(A,B)\dashv\vdash B.
\]
For this selector construction, first elaborate each $Q$ as a $D$ with
an unused witness. When the resulting formulas have unequal binder depths,
insert a vacuous witnessed gate to align them. The exact structural recursion and its
proof appear in the inherited native-basis appendix.

Let $\delta(v)=C(v,1,0)$ and $p(z)=I(z,0;u.\delta(u))$.
The following are formula macros, not additional primitive logical rules:
\begin{align*}
 \bot_* &:= (0=S0),\\
 A\Rightarrow_* B&:=Q_y M_y(A,B),\\
 A\wedge_*B&:=Q_y G_yM_{p(y)}(A,B),\\
 \exists_*x\,A(x)&:=D_{x,y}G_yA(x),\\
 A\vee_*B&:=\exists_*v\,M_v(A,B),\\
 \forall_*x\,A(x)&:=A(0)\wedge_*Q_yA(y).
\end{align*}
Every newly bound variable is fresh; normalize all numerical term
expressions in the macros.

For implication, case introduction supplies a closure from the assumed
zero selector to the successor selector, and QE at $1$ gives elimination.
For conjunction the gated base is true, and the instances $1,2$ give the
two operands because $p(1)=1$ and $p(2)=0$. Numerical cases give its
introduction. The existential macro uses its explicit witness and a true
gated base; its elimination uses DE and QE at $1$. Disjunction uses
selector witnesses $0,1$ and numerical cases. Universal introduction
and elimination are exactly the base instance together with its
conditional closure. These derivations supply all ordinary intuitionistic
natural-deduction rules.

Under $0=1$, equality transport through $C(v,s,t)$ gives $s=t$ for
arbitrary terms. External induction on formula syntax then gives
every native formula, so the designated contradiction supports explosion.
No excluded-middle schema for quantified formulas is introduced.

\subsection{The intermediate ordinary theory}

Let $T_I$ be intuitionistic first-order logic with actual equality and
genuine falsity, the scalar terms, their two computation equations,
contextual induction for literal equations, and
\[
 \forall x\,(x=0\vee\exists y\,x=Sy),
 \qquad (0=S0)\to\bot.
\]
The last sentence is an explicit falsity bridge. Positive equations alone
would permit a singleton structure, so this bridge cannot be silently
discarded from the ordinary realization.

The ordinary interpretation is
\[
 (s=t)^\circ=(s=t),\quad
 (Q_yB)^\circ=B^\circ(0)\to\forall y\,B^\circ(y),\quad
 (D_{x,y}B)^\circ=
 \exists x\bigl(B^\circ(x,0)\to\forall y\,B^\circ(x,y)\bigr).
\]
Case introduction is sound by coverage. Equation-progress introduction
is sound by equation induction under the temporarily assumed base.
The other rules are ordinary equality and logical reasoning.
Conversely the displayed native macros, equation induction derived
from Q-progress and QE, numerical cases, and explosion reconstruct
$T_I$. Coverage itself follows by cases, using the native disjunction
and existential macros.

The key native formula composite is
\[
 D_{x,y}B
 \dashv\vdash
 \exists_*x\bigl(B(x,0)\Rightarrow_*\forall_*y\,B(x,y)\bigr).
\]
Forwards, DE opens a witness; under its base the exported closure
gives every instance. Introduce the three macros and close DE.
Backwards, open the existential at a fresh witness, assume its base,
and use implication and universal elimination at $Sz$ to introduce
the original $D$ by cases. The analogous $Q$ composite is its conditional
closure. Induction on formulas proves both complete formula round trips.

\subsection{Arithmetic compilation and both numerical composites}

Each fixed scalar update is elaborated as a function symbol.
There is no variable ranging over function codes and no uniform
object-language evaluator. In one direction, iPRA defines a fixed
iterator by primitive recursion whose update ignores its recursion
counter. Its equations and every literal-equation induction instance
are then available.

For the reverse direction, the inherited internal scalar proofs proceed
in the following order:
\begin{enumerate}
\item Construct a scalar predecessor $P$ and prove $P0=0$ and $P(Sn)=n$.
This uses equation induction and numerical cases, without assuming
pairing, equality reflection, or successor injectivity.
\item Define truncated subtraction by iteration of $P$ and construct a
Boolean characteristic term for \emph{actual} equality. Prove the
reflection equivalence throughout the numerical domain.
\item Translate every fixed quantifier-free formula into its
characteristic equation. Contextual equation induction therefore
supplies every quantifier-free induction instance of iPRA.
\item Construct pairing and projections by bounded search, and compile
every fixed primitive-recursive definition into scalar iteration.
\end{enumerate}

For example, with pairing $\pi$ and projections $L,H$ already proved,
a primitive recursion
\[
 f(0,\vec p)=g(\vec p),\qquad
 f(Sn,\vec p)=h(n,f(n,\vec p),\vec p)
\]
is represented by
\[
\begin{aligned}
 V(n,\vec p)&=
 I\bigl(n,\pi(0,g^*(\vec p));
 u.\pi(SL(u),h^*(L(u),H(u),\vec p))\bigr),\\
 f^*(n,\vec p)&=H(V(n,\vec p)).
\end{aligned}
\]
Equation induction proves $L(V(n,\vec p))=n$; projection equations
then give both defining equations for $f^*$. In $\mathcal F$, use
the normal form of this finite term. It has the same proved
defining equations by the normalization comparison.

Uniqueness is internal, not an assertion about standard natural numbers.
For two iterators with the same initializer and pointwise equal updates,
induct on equality of their values. For ordinary primitive recursion,
use the same argument with its counter-dependent update. External
induction on a finite function description shows that its double
translation satisfies its original equations and hence is provably
equal to the original function.

Both numerical interpretations have domain formula $x=x$ and retain
actual equality. Consequently the identity graph $x=y$ is internally
total, functional, bijective, and preserves each interpreted operation.
Combined with the formula composites, this proves bi-interpretability.
Finite proof translations preserve $0=S0$ and ordinary falsity in both
directions, giving equiconsistency. This does not assert that either
theory proves its own consistency.

\paragraph{Status of the inherited formalization.}
The pinned repository supplies Lean statements for native/iPRA
faithfulness, both formula round trips, and the contradiction equivalence.
Those source endpoints were inspected; the upstream development was
not rebuilt here. The new full-normal-form calculus, its normalization
metatheory, and the Python correspondence have written proofs and
executable checks, but no new Lean or Rocq verification is claimed.
